We re-optimised \nRevAttackers{} held-out attackers against the revised stack:
three CEM seeds per (family, FPR) cell plus one tail-objective TPE. Each was
evaluated on \nROneEvalN{} missions under three wind conditions, from a seed
block never used before (\nRevUndet{} undetected-attack missions).
\begin{itemize}
\item \emph{No violation under either budget.} With the R1 budget, the test found
  \nRevEmpViol{} Holm-significant violations in \nRevEmpTests{} tests; with the
  LP budget, \nRevLpViol{} in \nRevLpTests. Superposition was exact on every
  undetected mission (largest excess \nRevSupUnsat\,m). Missions like the benign-tail failures that broke the pilot instantiation
  saturate, so they now raise the alarm and leave the event $U$.
\item \emph{The empirical budget was exceeded again.} In \nRevEmpExceed{} of
  \nRevEmpCells{} cells a held-out attacker went beyond $\dadv=\nROneDadvFive$\,m,
  up to \nRevAdvMax\,m (\nRevAdvMaxFam, \nRevAdvMaxOpt). Benign slack absorbed it
  and no violation followed, but the stealth budget we can use is still not a
  sound constant. The LP budget was never exceeded. It can be exceeded only with probability
  $\epsilon$, or on missions with a GNSS outage, which its proof does not cover.
\item \emph{What the bound certifies.} With the R1 budget, the floor at 5\% FPR
  is \nRevEmpFloorTen{} at $m=10$\,m and \nRevEmpFloorTwenty{} at $m=20$\,m.
  With the LP budget it is \nRevLpFloorTen{} at both, so it certifies nothing
  in the corridors we test.
\item \emph{Limited power.} Of the \nRevEmpInformative{} tests at which the bound
  is informative, \nRevEmpFailInf{} saw any undetected failure. The revision is
  therefore \emph{not refuted}, not \emph{confirmed tight}. The negative
  controls remain the evidence that the test can fail, and the scaled run of
  \cref{sec:scaled} shows that it does when a condition is missing.
\end{itemize}
